\section{Large Momentum Vs. Short Distance Expansion}

The Euclidean correlator in Eq. (5) introduced in Ref. \cite{Ji:2013dva}
has also been considered in coordinate-space factorizaton (CSF)~\cite{Radyushkin:2017cyf}, which was
introduced in an early work on meson DAs with current-current correlators
\cite{Braun:2007wv} (see also~\cite{Ma:2017pxb}).
The correlator can be factorized in terms of the LF correlations
with expansion parameter $(z\Lambda_{\rm QCD})^2$. The formalism is
naturally suited for calculating moments of PDFs or short-distance
LF correlations. To obtain the full parton physics, however, one has to
simultaneously consider the constraint on the external momentum $P^z\sim 1/z \gg \Lambda_{\rm QCD}$.
This is identical to the observation in Ref. \cite{Ji:2014gla}: One must use
large momenta to capture the full dynamical range of PDFs, which requires information on
long-range correlations in $\lambda$.
Despite their formal equivalence~\cite{Ji:2017rah,Izubuchi:2018srq,Radyushkin:2017lvu},
some analytical matching calculations might
more conveniently be done in coordinate space. Not surprisingly, the same LaMET
lattice data are needed for a CSF analysis to get the PDFs.
Nominally, CSF can also admit data at small $P^z$, but the same information is
contained already in large $P^z$ data at smaller $z$.

The CSF expansion is formulated in terms of
the Euclidean distance $z$, which is required to be small, i.e.,
\begin{equation}
         z\ll 1/\Lambda_{\rm QCD}\ ,
\end{equation}
to ensure the validity of perturbation theory and leading-twist dominance.
Assuming the largest $z$ for the leading-twist
approximation to be $z_{\rm LT}$ (say, the value of $z$ for which the higher-twist contribution
is at the level $\sim20\%$), then the small
expansion parameter is $(z_{\rm LT}\Lambda_{\rm QCD})^2\ll 1 $ when
potential linear divergences are subtracted
before the expansion is made. Therefore, only the matrix element
of $O(z)$ within the range $[0, z_{\rm LT}]$ has a simple interpretation
in terms of leading-twist parton physics.

An interesting question is then: What is the value of $z_{\rm LT}$?
If we take $\Lambda_{\rm QCD}\sim 300$ MeV, and  $z_{\rm LT}\Lambda_{\rm QCD}=1/2\sim 1/3$ as
a small parameter, then the estimate is that $z_{\rm LT}$ is around $0.25 \sim 0.33$ fm. An
upper limit is probably 0.4 fm. A good estimate of $z_{\rm LT}$ can be provided by comparing the matrix element $\langle P=0| O(z)|P=0\rangle$ or $\langle \Omega| O(z)|\Omega\rangle$, both of which have been proposed to renormalize the bare quasi-LF correlation~\cite{Radyushkin:2017cyf,Braun:2018brg,Li:2020xml}, to the leading-twist contributions in their OPE.

Let us take the zero-momentum matrix element  for the isovector case as an example. In the $\MS$ scheme, it has a short distance expansion of the form~\cite{Radyushkin:2017lvu,Izubuchi:2018srq}
\begin{align}
			\tilde{h}(z,\mu,P^z\!=\!0)&= {1\over 2M}\langle P=0| \bar{\psi}(z)\gamma^0 W(z,0)\psi(0)|P=0\rangle \nn\\
			& = c_0(\mu^2 z^2) a_0 + {\cal O}(z^2\Lambda_{\rm QCD}^2)\,,
\end{align}
where $W(z,0)$ is a spacelike straight gauge link. Here $\mu$ is the $\MS$ renormalization scale, and $a_0=1$ is the conserved lowest moment of the correpsonding twist-2 PDF. The one-loop Wilson coefficient $c_0=1/  Z^{\overline{\rm MS}}_{\rm ratio}$ with $Z^{\overline{\rm MS}}_{\rm ratio}$ given in \eq{zmsr}~\cite{Izubuchi:2018srq},
and the two-loop result can be found in Ref.~\cite{Li:2020xml}.

According to \eq{ren}, the mass-subtracted matrix element includes logarithmic divergences that are independent of $z$ and should not constitute significant corrections in lattice perturbation theory. Therefore, we can roughly approximate its OPE by replacing $\mu$ with $1/a$,
\begin{align}\label{eq:subtracted}
	e^{-\delta m |z|}\tilde{f}(z,a,P^z\!=\!0) \!=\! c_0(z^2/a^2) \!+\! {\cal O}(z^2\Lambda_{\rm QCD}^2, a^2\!/\!z^2)\,,
\end{align}
where the lattice discretization effects are expected to be of ${\cal O}(a^2/z^2)$.
Since the lattice matrix elements are convergent as $z\to0$, which is contrary to the logarithmically divergent behavior in the $\MS$ OPE, we expect the above approximation to be reliable within the range $a < |z| < z_{\rm LT}$ where the discretization and higher-twist effects are both suppressed. Note that lattice OPE is usually complicated by the broken Lorentz symmetry and operator mixings. Nevertheless, since for $P^z=0$ the only leading-twist contribution comes from the conserved vector current, we can ignore such effects here.

In \fig{ope} we plot the mass-renormalized pion lattice matrix element. The bare lattice matrix element comes from a recent calculation of the pion valence PDF on an ensemble with $a=0.06$ fm and pion mass $m_\pi=300$ MeV~\cite{Gao:2020ito}. On the same lattice ensemble, the Wilson-line mass correction $\delta m$ was fitted from the quark-antiquark potential~\cite{Izubuchi:2019lyk}, and its value is given in lattice units as $a \delta m = -0.1568$. The leading-twist contribution is plotted with next-to-leading-order (NLO) corrections and NLO correction plus LL resummation for fixed $\alpha_s$,
\begin{align}
	\left[ 1 + {\alpha_s(1/a)C_F\over 2\pi}  {5\over2}\right]\left({z^2\over 4e^{-2\gamma_E}a^2}\right)^{{\alpha_s(1/a)C_F\over 2\pi}{3\over 2}}\,.
\end{align}
Here we choose $\alpha_s$ in the $\MS$ scheme at scale $1/a$ as the input for OPE, which should allow for better convergence than the bare lattice coupling~\cite{Lepage:1992xa}. To estimate the uncertainty from the choice of $\alpha_s$, we vary the $\MS$ scale from $1/(2a)$ to $2/a$. Though a standard procedure of improvement shall be performed to define $\alpha_s$ on the lattice, we expect that it will not alter the following conclusion.

As one can see, for $z\le a$, the lattice result is significantly different from the leading-twist approximations due to discretization effects. As $z$ increases, the agreement becomes better. However, for $z\ge 0.3$ fm, the lattice result starts to deviate dramatically from the leading-twist approximations, showing that the higher-twist contributions become significant. Therefore, we can roughly estimate that $z_{\rm LT}\sim 0.3$ fm.

One can also look at the case of the better established heavy-quark potential. It is well-known that the static heavy-quark potential receives both perturbative and non-perturbative contributions. The  perturbative static potential is known up to ${\rm N^3LO}$ level~\cite{Smirnov:2009fh}
and can be expressed in terms of the QCD running coupling constant. In Refs.~\cite{Bazavov:2012ka,Bazavov:2014soa,Bazavov:2019qoo}, the running coupling constant has been extracted from lattice calculation of the static energy at short distances. The ${\rm N^3LL}$ perturbative result agrees well with lattice data up to $r \sim 0.2$ ${\rm fm}$. However, it is well-known that the perturbative series for the static potential suffers from a renormalon ambiguity~\cite{Beneke:1998rk,Hoang:1998nz} and breaks down at large distance. The non-perturbative heavy-quark potential can be simulated using lattice QCD, and is well-known to be dominated by the linear term of the form $\sigma r$ at large distance. Phenomenologically, the static potential can be well approximated by the linear+Coulumb QCD static potential, $V(r)=-e/r+\sigma r$ where $e= 0.25\sim 0.5$ while $\sqrt{\sigma} \approx 477$ ${\rm MeV}$~\cite{Eichten:1974af,Bali:2000gf,Aubin:2004wf}. When the perturbation theory is about to break down, the perturbative contribution $-e/r$ and the confining contribution $\sigma r$ should be of the same order of magnitude, which determines $r_c \approx\sqrt{e/\sigma}$ to be around $0.2\sim 0.3$ ${\rm fm}$. This is consistent with the result in Refs.~\cite{Bali:1999ai,Necco:2001gh,Bazavov:2012ka}. The boundary $z_{\rm LT}$ should be of the same order of magnitude.

\begin{figure}[t]
\centering
\includegraphics[width=1\columnwidth]{images/ZLT_pion}
\caption{Comparison of the mass-renormalized pion lattice matrix element~\cite{Gao:2020ito} and its leading-twist approximation with NLO and NLO$+$LL corrections for fixed $\alpha_s$, respectively.
The strong coupling is $\alpha_s(1/a)=0.242$, and the error band was obtained by varying $\alpha_s$ from $\alpha_s(1/(2a)$ to $\alpha_s(2/a)$. This shows that the leading-twist approximation becomes unreliable at $z_{\rm LT}\sim 0.3$ fm.}
\label{fig:ope}
\end{figure}

\begin{figure}[t]
\includegraphics[width=1\columnwidth]{images/potential}
\caption{A comparison between the perturbative potential in $\alpha_s^3$ order to the linear+Coulumb QCD static potential taken from~\cite{Bali:2000gf,Aubin:2004wf}. The shaded area corresponds to the perturbative prediction with the renormalization scale $\mu$ ranging from $r/2$ to $3r/2$. The black line corresponds to $\mu=1/r$. From the figure, it is clear that beyond $0.2$ to $0.3$ ${\rm fm}$, the perturbative potential starts to deviate from the full non-perturbative results and becomes unreliable.}
\label{fig:poten}
\end{figure}

With the estimated $z_{\rm LT}$ above, one can also define
\begin{equation}
       \lambda_{\rm LT} = P^zz_{\rm LT},
\end{equation}
then the matrix element for the quasi-LF distance $[0, \lambda_{\rm LT}]$
can be used to extract parton distributions with the matching formula in \eq{coordfact}~\cite{Radyushkin:2017cyf,Ji:2017rah,Izubuchi:2018srq}.
Thus, the coordinate-space approach is useful for extracting the LF correlation functions in a limited range, with the LF distance ranging between $[0,\lambda_{\rm LT}]$.

The CSF approach has also been used for products of currents made of
quark bilinears~\cite{Detmold:2005gg,Braun:2007wv,Bali:2017gfr,Bali:2018spj,Ma:2017pxb,Detmold:2018kwu,Sufian:2019bol,Sufian:2020vzb}.
Renormalization of power divergences in the
CSF expansion for the quark and gluon blinears with Wilson line
is easier to handle. In particular, a version of the ratio method
which  divides by the matrix element at zero momentum, can be used to
eliminate the power divergences in the lattice matrix elements~\cite{Radyushkin:2017cyf,Orginos:2017kos}.
On the other hand, the current products can also be used
in LaMET expansion after Fourier transforming into momentum
space~\cite{Ji:2020ect}.

However, the CSF does not allow for directly calculating the $x$-distribution, because one needs
the LF correlation at all LF distances. The requirement for $z\le z_{\rm LT}$ makes it unfeasible
to reach large $\lambda$ values for the Fourier transform with the largest momentum on contemporary lattice resources. Thus, to reconstruct
the PDF, one has to parameterize the functional form of the $x$ dependence, just like that
in the phenomenological fits, and then
inverse Fourier transform it to the coordinate space and fit to a limited range
of LF correlations~\cite{Orginos:2017kos,Karpie:2018zaz,Joo:2019bzr,Joo:2019jct,Sufian:2019bol,Joo:2020spy,Sufian:2020vzb,Bhat:2020ktg,Gao:2020ito,Fan:2020cpa}.
This process is hardly under control, because it is difficult to estimate
the systematic uncertainty from the parameterization or assumptions of the PDF. As evident in the fits performed in literature so far~\cite{Joo:2019bzr,Joo:2019jct,Sufian:2019bol,Joo:2020spy,Gao:2020ito,Fan:2020cpa}, either the errors in the end-point regions become smaller and smaller, or the errors shrink to almost zero for certain moderate values of $x$. These imply unaccounted systematics from the artifacts of the particular model used, which is also reflected by their inconsistency with global fits that use similar parameterizations.

From a different angle, the above practice amounts
to postulating (or modeling) certain correlation between short- and long-
distance behaviors of the LF correlations. Such a postulation
has no first-principles foundation, and it can happen that the lattice data in
the limited range of $\lambda$ be fitted equally well by more than one parameterizations which
have completely different asymptotic behaviors~\cite{Sufian:2020vzb}.

Despite the difficulty in providing a controlled calculation of the $x$-distribution, the CSF method allows for model-independent extraction of the Mellin moments using the OPE. Nevertheless, with limited range of $\lambda$, the LF correlations will be sensitive to only the lowest ones, which is related, but not in a direct one-to-one correspondence, to the predictive power in the $x$-space.

To make a more direct comparison between the momentum space and coordinate space approaches, let us consider the following example. Assume that the quasi-LF correlations defining the quasi-PDFs behave like
\begin{align}
\tilde h(\lambda,z)=h(\lambda)e^{-m|z|} \ ,
\end{align}
with $h(\lambda)$ being the light-cone correlator. The exponential $e^{-m|z|}$ with $m\sim \Lambda_{\rm QCD}$ is used to model higher-twist contributions.
From this equation, it is clear that if one stays in position space, the CSF is accurate only when $z\ll 1/m$. The available range of $\lambda$ is therefore much smaller than $P^z/m$, which indicates that the number of moments one can access is much less than $P^z/m$.

From the discussion above, it is clear that the momentum and coordinate space expansions are different expansion schemes. Even though they are equivalent in the infinite momentum limit, they are different at finite
momentum $P^z$.
In the latter, the information is filtered directly
in coordinate space. One gets parton correlations
in a finite range of LF distance which correspond to
the number of moments controlled by $1/P^z$. In contrast, the former uses all the coordinate space information, filtering
higher-twist physics in momentum space through $1/(y^2(P^z)^2)\ll 1$ and $1/((1-y)^2(P^z)^2)\ll 1$.
Therefore, one gets partron distributions in an interval of $x$
with systematic control of errors, which can be directly compared
with experimental data.

Finally, we remark that the relative size of the perturbative correction to the quasi-PDF depends on $x$, where one usually observes larger corrections in the small- and large-$x$ regions. On the other hand, although the size of the perturbative corrections in the coordinate space is usually small for finite $\lambda$, it can still lead to significant corrections in the end-point regions in momentum space.
